% TOP_PROBLEM: {"id":30007694,"problem_number":"LOCAL-30007694","title":"Democratizing innovation","category_id":21,"category":{"id":21,"name":"economics","display_name":"Economics","description":"Open economic theory statements, published research agendas, and proposed scoped specifications; question provenance and historical priority are qualified per record.","slug":"economics","order_index":21,"created_at":"2026-10-08T00:00:00Z"},"status":"research_question","external_url":null,"legacy_ids":[],"published":false,"statement":"Measure whether shared AI research access narrows the innovation gap between eligible small firms and incumbents while increasing aggregate net value.","economics":{"original_id":183,"field":"Innovation and AI economics","kind":"empirical","formalization_basis":"adapted_research_question","source_status":"explicit_open","priority_status":"earliest_found","priority_note":"Earliest relevant explicit question or agenda passage located in this audit. No claim of first-ever formulation; earlier versions and later solutions were not exhaustively traced.","earliest_verified_reference":{"authors":["Erik Brynjolfsson","Anton Korinek","Ajay K. Agrawal"],"title":"A Research Agenda for the Economic Implications of Transformative Artificial Intelligence","year":2024,"url":"https://digitaleconomy.stanford.edu/app/uploads/2024/11/ETAI-White-Paper.pdf","doi":null,"locator":"Section 3.2, printed pp. 9–10, democratization of discovery question","evidence_summary":"Explicitly asks whether AI democratizes innovative capacity across organizations."},"current_reference":{"authors":["Erik Brynjolfsson","Anton Korinek","Ajay K. Agrawal"],"title":"A Research Agenda for the Economic Implications of Transformative Artificial Intelligence","year":2024,"url":"https://digitaleconomy.stanford.edu/app/uploads/2024/11/ETAI-White-Paper.pdf","doi":null,"locator":"Section 3.2, printed pp. 9–10, democratization of discovery question","evidence_summary":"Explicitly asks whether AI democratizes innovative capacity across organizations."},"formalization_note":"The source verifies a broader question or research agenda; the population, estimand, equations and restrictions here are our scoped adaptation. This is not a quoted canonical conjecture, and the audit does not prove contemporary unsolved status for every special case."},"statusQualification":"Published question or research agenda verified at the cited locator. The mathematical specification is adapted for this catalogue and includes additional assumptions; source publication does not establish first-ever priority or certify this exact model remains unsolved. The source verifies a broader question or research agenda; the population, estimand, equations and restrictions here are our scoped adaptation. This is not a quoted canonical conjecture, and the audit does not prove contemporary unsolved status for every special case.","statusReviewedAt":"2026-10-08"}
% FIRST_POSTED: 2026-10-08
% ENTRY_KIND: statement_only
% !TeX program = lualatex
\documentclass[11pt]{article}
\usepackage{fontspec}
\usepackage[a4paper,margin=25mm]{geometry}
\usepackage{amsmath,amssymb,mathtools}
\usepackage{xurl}
\usepackage[hidelinks]{hyperref}
\newcommand{\sref}[2]{\hyperlink{source:#1:#2}{[S#2]}}
\setlength{\emergencystretch}{3em}
\begin{document}
% problemId:30007694
\section*{Democratizing innovation}\label{30007694}
\subsection{Definitions and mathematical statement}
Consider firms eligible for a public AI research-access program, partitioned before treatment into small firms \(G=s\) and resource-rich incumbents \(G=l\) using a stated employment threshold. Let \(A=1\) provide an identical subsidized model, computation allowance and training package, while \(A=0\) retains existing access. Over three years, define \(Y_i(a)\) as the value of independently validated new products from firm \(i\), with zero for nonparticipation or unsuccessful research. Group effects are
\[
\tau_g=E[Y_i(1)-Y_i(0)\mid G=g],\qquad g\in\{s,l\},
\]
and the change in the average innovation gap is \(\Gamma=\tau_s-\tau_l\). This gap does not itself measure aggregate welfare, so separately record entry, exit, innovation expenditure and product diffusion. Use program eligibility lotteries or clustered rollout with recorded access compliance; account for knowledge spillovers and changes in incumbent competition within product markets. Define a market-level policy effect by assigning access to all eligible firms in a market and measuring total validated value net of program resource costs. Which access and support arrangements make \(\Gamma>0\) while increasing this aggregate outcome? The scope is eligible firms in the selected sectors, with uncertainty about commercialization and crowding out explicitly measured rather than an unrestricted claim that AI democratizes innovation.

\par\textbf{Formulation and scope.} The source verifies a broader question or research agenda; the population, estimand, equations and restrictions here are our scoped adaptation. This is not a quoted canonical conjecture, and the audit does not prove contemporary unsolved status for every special case.

\par\textbf{Open-question provenance.} An explicit question or research agenda was verified at the source locator. This does not by itself certify that every later version remains unresolved \sref{30007694}{1}.

\par\textbf{Historical priority.} Earliest relevant explicit question or agenda passage located in this audit. No claim of first-ever formulation; earlier versions and later solutions were not exhaustively traced.

\subsection{Short English statement}
Measure whether shared AI research access narrows the innovation gap between eligible small firms and incumbents while increasing aggregate net value.

\subsection{Sources}
\begin{itemize}\raggedright
\item[\textnormal{[S1]}]\hypertarget{source:30007694:1}{} Erik Brynjolfsson, Anton Korinek, Ajay K. Agrawal. 2024. A Research Agenda for the Economic Implications of Transformative Artificial Intelligence. Section 3.2, printed pp. 9–10, democratization of discovery question.
\url{https://digitaleconomy.stanford.edu/app/uploads/2024/11/ETAI-White-Paper.pdf}.
{\small\textit{Earliest verified question/agenda reference; historical priority qualified below. Explicitly asks whether AI democratizes innovative capacity across organizations.}}
\end{itemize}
\end{document}
